\section{因果结构的代数形式}
\label{sec:12-Real-Analysis-Support/10-Essential-Structures/04}

一家医院的记录显示，用过某型号呼吸机的病人，院内死亡率明显高于没用的。采购部门要据此停止续购，主治医生的解释却是：能上这台机器的都是最重的病人。两种说法对同一份数据给出一模一样的相关性，对``再买几台会怎样''却给出相反的建议。手上的数据能不能分辨它们？

这一节的回答是不能，而且不是数据量的问题。

\subsection{同一个联合分布可以由不同机制产生}

先把机制写清楚。结构因果模型由一组赋值构成，
\[
X_j := f_j\big(\mathrm{pa}_j,\,U_j\big),\qquad j=1,\dots,d,
\]
其中 \(\mathrm{pa}_j\) 是 \(X_j\) 的父节点集，\(U_j\) 是外生噪声，符号 \(:=\) 表示赋值。它与等号的差别是代数上的：等式 \(Y=aX+U\) 可以移项写成 \(X=(Y-U)/a\)，赋值不行。父子关系连成一张有向无环图 \(G\)，这张图就是机制的记账方式。

现在看一个数字例子。取二值 \(U\)，\(P(U=1)=0.5\)，并设
\[
P(X=1\given U=u)=P(Y=1\given U=u)=
\begin{cases}0.9,&u=1,\\ 0.1,&u=0,\end{cases}
\]
两条赋值中 \(X\) 与 \(Y\) 各用独立噪声。算出来的联合分布满足 \(P(X=1)=P(Y=1)=0.5\)，\(P(X=1,Y=1)=0.41\)，于是
\[
P(Y=1\given X=1)=0.82,\qquad P(Y=1\given X=0)=0.18 .
\]
观察到的差是 \(0.64\)。但这个模型里 \(X\) 对 \(Y\) 没有任何箭头，强行设定 \(X\) 的值不会改变 \(Y\) 的分布，干预效应是 \(0\)。

同一份联合分布也能由 \(X\to Y\) 生成：令 \(P(X=1)=0.5\)，\(P(Y=1\given X=1)=0.82\)，\(P(Y=1\given X=0)=0.18\)。逐格核对，两张表完全相同。可这个模型里干预效应就是 \(0.64\)。再把箭头掉过来，令 \(P(Y=1)=0.5\)、\(P(X=1\given Y=1)=0.82\)、\(P(X=1\given Y=0)=0.18\)，联合分布仍然一样，而干预 \(X\) 对 \(Y\) 毫无影响。

\begin{center}
\begin{tikzpicture}[>=stealth,line width=0.8pt]
\node at (4.9,2.6) {\footnotesize 三张图给出同一个 \(P(X,Y)\)};
\node[draw,circle,minimum size=7mm,inner sep=1pt] (ax) at (0,0.9) {\(X\)};
\node[draw,circle,minimum size=7mm,inner sep=1pt] (ay) at (1.8,0.9) {\(Y\)};
\draw[->] (ax) -- (ay);
\node at (0.9,-0.15) {\footnotesize (a) 效应 \(0.64\)};
\node[draw,circle,minimum size=7mm,inner sep=1pt] (bx) at (3.9,0.9) {\(X\)};
\node[draw,circle,minimum size=7mm,inner sep=1pt] (by) at (5.7,0.9) {\(Y\)};
\draw[->] (by) -- (bx);
\node at (4.8,-0.15) {\footnotesize (b) 效应 \(0\)};
\node[draw,circle,minimum size=7mm,inner sep=1pt] (cx) at (7.8,0.9) {\(X\)};
\node[draw,circle,minimum size=7mm,inner sep=1pt] (cy) at (10.0,0.9) {\(Y\)};
\node[draw,circle,dashed,minimum size=7mm,inner sep=1pt] (cu) at (8.9,2.0) {\(U\)};
\draw[->] (cu) -- (cx);
\draw[->] (cu) -- (cy);
\node at (8.9,-0.15) {\footnotesize (c) 效应 \(0\)};
\end{tikzpicture}
\end{center}

\emph{观察数据把三张图放在一起，干预把它们分开：同一列数字，三个不同的采购建议。}

\subsection{条件独立能从图上读出来，方向读不出来}

图与分布之间的桥是 d-分离。给定 \(G\)，若节点集 \(A\) 与 \(B\) 被 \(C\) d-分离，则对每个按 \(G\) 分解的分布都有 \(A\perp B\given C\)。这是定理，方向是从图到分布。反向的``图上不分离就一定不独立''叫忠实性，它是一条附加假设，且可以被参数的精确抵消破坏——两条路径的效应恰好相消时，独立性出现在分布里却读不出图上。这层区分在\hyperref[sec:05-Probability/02-Conditional-Probability/04]{``条件独立与信息结构''}一节里有更细的讨论，那里也解释了为什么碰撞点 \(X\to C\leftarrow Y\) 上条件化会\emph{制造}依赖，而混杂 \(X\leftarrow Z\to Y\) 上条件化会消去依赖。两种操作外形一样，后果相反，\hyperref[sec:13-Machine-Learning/12-Causal-Inference/03]{``混杂、碰撞点与后门调整''}一节按图给出了该调整谁的判据。

那么观察数据最多能定到哪一步？答案是精确的。

\begin{theorem}
（Verma--Pearl）两张有向无环图蕴涵完全相同的条件独立关系，当且仅当它们有相同的骨架与相同的 v-结构。
\end{theorem}

上图里 (a) 与 (b) 骨架相同、都没有 v-结构，因此 Markov 等价；它们蕴涵的条件独立集合都是空集，任何数据都区分不了。这不是样本量不够，是等价类内部本来就没有可供区分的观察后果。

\subsection{干预删掉一整个因子，条件化只做重新归一}

把干预写成代数，差别就一目了然。按 \(G\) 分解的联合分布是 \(\prod_j P(x_j\given \mathrm{pa}_j)\)。做 \(\operatorname{do}(X_k=x_k^*)\)，等于把第 \(k\) 条赋值整条换成常数，图上剪掉所有指向 \(X_k\) 的箭头，于是
\[
P\big(x_1,\dots,x_d\given \operatorname{do}(X_k=x_k^*)\big)=
\begin{cases}
\displaystyle\prod_{j\neq k}P(x_j\given \mathrm{pa}_j), & x_k=x_k^*,\\[2mm]
0,&\text{否则.}
\end{cases}
\]
对照条件化：\(P(x_1,\dots,x_d\given X_k=x_k^*)\) 保留全部 \(d\) 个因子，只除以一个 \(P(X_k=x_k^*)\)。一个抹掉因子，一个重新归一，代数形式不同，数值自然不同。两者相等的充分条件是 \(X_k\) 没有父节点——随机化试验做的正是这件事，它用抛硬币替换掉原来的赋值。\hyperref[sec:13-Machine-Learning/12-Causal-Inference/02]{``结构因果模型把干预写成机制替换''}一节把这套记账展开到多变量。

从截断分解可以直接推出后门公式：若 \(Z\) 中没有 \(X\) 的后代，且 \(Z\) 阻断了每一条指向 \(X\) 的路径，则
\[
P(y\given \operatorname{do}(X=x))=\sum_z P(y\given x,z)\,P(z).
\]
右端全是观察量，所以效应可识别。这里必须把话说满：公式是定理，前提是那张图正确。图从哪来，公式不管。

\subsection{要越过等价类，必须付出假设}

观察数据决定联合分布，联合分布决定条件独立关系，条件独立关系（在忠实性下）决定 Markov 等价类。这条链走到头就停了。想再往前，得从别处拿东西。

随机化是最干净的一种：它不是从数据里推出因果，而是动手把 \(X\) 的父节点清空，让观察分布与干预分布重合。工具变量换一种付法，它要求存在 \(W\) 只通过 \(X\) 影响 \(Y\)，这条排除性约束本身不可由数据检验。函数形式也是一种通货：二元线性高斯模型的方向在观察上不可识别，而 Shimizu 等人的 LiNGAM 结果表明，把噪声换成非高斯之后，线性无环模型的整张图可由观察数据唯一确定。识别性由此从无到有，代价是``噪声非高斯且关系线性''这条假设，它只能被部分检验。

未观测混杂则是另一类。图 (c) 里的 \(U\) 若无法测量，任何观察数据都补不上这个缺口，可做的是敏感性分析：假定未观测混杂的强度不超过某个界，反推效应估计能被推翻到什么程度。得到的是区间与断点，不是点估计。\hyperref[ch:120]{``因果推断：预测世界不等于改变世界''}一章按这条线索排布了识别、估计与外推三层问题。

\subsection{一份因果结论要交代三样东西}

回到开头的呼吸机。图 (a) 的解读说机器有害，图 (c) 的解读说病情既决定上机也决定结局，两者对同一列数字给出的干预效应差了 \(0.64\)。要在它们之间做选择，数据不够，需要的是关于机制的陈述：谁决定了上机，有没有被记录下来。

因此一份因果结论应当交代清楚三样东西——所设的图或等价的排除性约束，识别公式如何从该图推出，以及在假设被推翻时结论会移动多少。前两样属于定理层面，可以逐步核验；第三样是稳健性，只能给界。缺了任何一样，剩下的就退回为相关性陈述，它在预测下一个观察值时依然有用，在回答``要不要改变什么''时无效。
